% 共用版式。xelatex 编译；Overleaf 上把 Songti SC / Heiti SC / Menlo 换成 FandolSong / FandolHei / DejaVu Sans Mono
\usepackage{xeCJK}
\setCJKmainfont{Songti SC}[BoldFont=Heiti SC]
\setCJKsansfont{Heiti SC}
\setCJKmonofont{Heiti SC}
\setmonofont{Menlo}[Scale=0.88]
\usepackage[a4paper,margin=1.8cm]{geometry}
\usepackage{amsmath}
% 数学也用 OpenType 字体：老式 Type 1 数学字体（CMSY / CMEX）在部分打印驱动上会丢「(」「−」
\usepackage{unicode-math}
\setmathfont{Latin Modern Math}
% unicode-math 下没有 \leadsto；LM Math 没有 \setminus 的字形（U+29F5），换成 \smallsetminus；\checkmark 正文里用，正文字体没有 ✓，借数学字体的
\newfontfamily\lmmathglyph{Latin Modern Math}
\AtBeginDocument{\let\setminus\smallsetminus\let\leadsto\rightsquigarrow\renewcommand{\checkmark}{{\lmmathglyph ✓}}}
\usepackage{booktabs,array,tabularx,longtable}
\usepackage{graphicx}
\usepackage[table]{xcolor}
\usepackage{listings}
\usepackage{enumitem}
\usepackage{titlesec}
\usepackage{fancyvrb}
\usepackage[most]{tcolorbox}
\usepackage[hidelinks]{hyperref}

\definecolor{codefg}{RGB}{150,30,60}
\definecolor{codebg}{gray}{0.95}
\definecolor{note}{gray}{0.38}
\definecolor{accent}{RGB}{40,80,150}

% 行内代码
\newcommand{\code}[1]{\texttt{\textcolor{codefg}{#1}}}
% 题号 + 题名
\newcommand{\lc}[2]{\textbf{#1}\,{\small #2}}

% Python 代码块：\begin{lstlisting} ... \end{lstlisting}
\lstset{language=Python,basicstyle=\ttfamily\small,backgroundcolor=\color{codebg},
  frame=none,xleftmargin=4pt,aboveskip=5pt,belowskip=5pt,columns=fullflexible,
  keywordstyle=\color{blue!60!black},commentstyle=\color{note},showstringspaces=false,
  breaklines=true,extendedchars=false}

% 可迁移的原则：\begin{principle} ... \end{principle}
\newtcolorbox{principle}{colback=accent!6,colframe=accent!60,boxrule=0.4pt,arc=1pt,
  left=6pt,right=6pt,top=3pt,bottom=3pt,fontupper=\small}
% 总览图（ASCII 图）：\begin{overview} ... \end{overview}
\DefineVerbatimEnvironment{overview}{Verbatim}{fontsize=\small,frame=single,framesep=6pt}
% 「你踩过 / 改过的地方」：\begin{mynote} ... \end{mynote}
\newtcolorbox{mynote}{colback=gray!5,colframe=gray!50,boxrule=0.4pt,arc=1pt,
  left=6pt,right=6pt,top=3pt,bottom=3pt,fontupper=\small,title={你踩过的坑},
  coltitle=black,colbacktitle=gray!15,fonttitle=\small\bfseries}

\titleformat{\section}{\Large\bfseries}{第\thesection 章}{0.8em}{}
\titleformat{\subsection}{\large\bfseries}{\thesubsection}{0.6em}{}
\titleformat{\subsubsection}{\normalsize\bfseries}{}{0em}{}
\titlespacing*{\section}{0pt}{0pt}{10pt}
\titlespacing*{\subsection}{0pt}{12pt}{4pt}
\titlespacing*{\subsubsection}{0pt}{8pt}{2pt}
\setlist{nosep,leftmargin=1.3em}
\renewcommand{\arraystretch}{1.3}
\setlength{\parindent}{0pt}
\setlength{\parskip}{4pt}
\renewcommand{\contentsname}{目录}
